
\documentclass{article}
\usepackage{colortbl}
\usepackage{makecell}
\usepackage{multirow}
\usepackage{supertabular}

\begin{document}

\newcounter{utterance}

\centering \large Interaction Transcript for game `eqbench', experiment `eq\_bench', episode 110 with reflexion{-}llama{-}3.1{-}8b{-}instruct.
\vspace{24pt}

{ \footnotesize  \setcounter{utterance}{1}
\setlength{\tabcolsep}{0pt}
\begin{supertabular}{c@{$\;$}|p{.15\linewidth}@{}p{.15\linewidth}p{.15\linewidth}p{.15\linewidth}p{.15\linewidth}p{.15\linewidth}}
    \# & \multicolumn{2}{c}{Player} && \multicolumn{2}{c}{Game Master} \\
    \hline

    \theutterance \stepcounter{utterance}  
    & & & \multicolumn{4}{p{0.6\linewidth}}{
        \cellcolor[rgb]{0.9,0.9,0.9}{
            \makecell[{{p{\linewidth}}}]{
                \texttt{\tiny{[P1$\langle$GM]}}
                \\ 
\texttt{Your task is to predict the likely emotional responses of a character in this dialogue:} \\
\\ 
\texttt{Luca: So, what do you think?} \\
\texttt{Charlie: It's... interesting.} \\
\texttt{Luca: Interesting? That's it?} \\
\texttt{Charlie: Well, it's not exactly a compliment, is it?} \\
\texttt{Luca: You haven't even tasted it yet. } \\
\texttt{Charlie: I don't need to taste it to know it's a mess.} \\
\texttt{Luca: A mess? You're calling my masterpiece a mess?} \\
\texttt{Charlie: If the shoe fits...} \\
\texttt{Luca: You're just a critic, Charlie. You don't know what it takes to create something.} \\
\texttt{Charlie: And you don't know what it takes to appreciate something. } \\
\texttt{Luca: I've been a chef for over 20 years. I've earned my stars.} \\
\texttt{Charlie: And I've been a critic for 15. I've earned my stripes.} \\
\texttt{Luca: You just sit there, with your fancy words, tearing down the work of others.} \\
\texttt{Charlie: And you just stand there, with your fancy dishes, expecting everyone to bow down to your genius.} \\
\texttt{Luca: You're just bitter because your life is falling apart.} \\
\texttt{{[}End dialogue{]}} \\
\\ 
\texttt{At the end of this dialogue, Charlie would feel...} \\
\texttt{Embarrassed} \\
\texttt{Provoked} \\
\texttt{Indifferent} \\
\texttt{Grateful} \\
\\ 
\texttt{Give each of these possible emotions a score from 0{-}10 for the relative intensity that they are likely to be feeling each.} \\
\\ 
\texttt{You must output in the following format, including headings (of course, you should give your own scores), with no additional commentary:} \\
\\ 
\texttt{Embarrassed: <score>} \\
\texttt{Provoked: <score>} \\
\texttt{Indifferent: <score>} \\
\texttt{Grateful: <score>} \\
\\ 
\\ 
\texttt{{[}End of answer{]}} \\
\\ 
\texttt{Remember: zero is a valid score, meaning they are likely not feeling that emotion. You must score at least one emotion > 0.} \\
\\ 
\texttt{Your answer:} \\
            }
        }
    }
     \\ \\

    \theutterance \stepcounter{utterance}  
    & \multicolumn{4}{p{0.6\linewidth}}{
        \cellcolor[rgb]{0.561,0.737,0.561}{
            \makecell[{{p{\linewidth}}}]{
                \texttt{\tiny{[P1$\rangle$GM]}}
                \texttt{Embarrassed: 0} \\
\texttt{Provoked: 9} \\
\texttt{Indifferent: 0} \\
\texttt{Grateful: 0} \\
            }
        }
    }
    & & \\ \\

    \theutterance \stepcounter{utterance}  
    & & & \multicolumn{2}{p{0.3\linewidth}}{
        \cellcolor[rgb]{0.9,0.9,0.9}{
            \makecell[{{p{\linewidth}}}]{
                \texttt{\tiny{[GM$|$GM]}}
                \texttt{Embarrassed: 0} \\
\texttt{Provoked: 9} \\
\texttt{Indifferent: 0} \\
\texttt{Grateful: 0} \\
            }
        }
    }
    & & \\ \\

    \theutterance \stepcounter{utterance}  
    & & & \multicolumn{2}{p{0.3\linewidth}}{
        \cellcolor[rgb]{0.9,0.9,0.9}{
            \makecell[{{p{\linewidth}}}]{
                \texttt{\tiny{[GM$|$GM]}}
                \texttt{Embarrassed: 2} \\
\texttt{Provoked: 8} \\
\texttt{Indifferent: 0} \\
\texttt{Grateful: 0} \\
            }
        }
    }
    & & \\ \\

    \theutterance \stepcounter{utterance}  
    & & & \multicolumn{2}{p{0.3\linewidth}}{
        \cellcolor[rgb]{0.9,0.9,0.9}{
            \makecell[{{p{\linewidth}}}]{
                \texttt{\tiny{[GM$|$GM]}}
                \texttt{game\_result = LOSE} \\
            }
        }
    }
    & & \\ \\

    \theutterance \stepcounter{utterance}  
    & & \multicolumn{4}{p{0.6\linewidth}}{
        \cellcolor[rgb]{0.95,0.95,0.95}{
            \makecell[{{p{\linewidth}}}]{
                \texttt{\tiny{[P1$\leftrightarrow$P]}}
                \texttt{Post{-}episode Reflexion (not sent to GM):} \\
\texttt{New plan: } \\
\\ 
\texttt{Upon re{-}examining the dialogue, I realized that I focused too much on Charlie's responses to Luca's statements, but didn't consider the overall tone and progression of the conversation. I should have taken into account the escalating nature of the argument, with both characters becoming increasingly defensive and critical of each other.} \\
\\ 
\texttt{To improve my strategy, I will pay closer attention to the emotional cues and body language implied in the dialogue, such as Luca's defensiveness and Charlie's sarcasm. I will also consider the context of the conversation, including the fact that Charlie is a critic and Luca is a chef, and how this power dynamic may influence their interactions.} \\
\\ 
\texttt{In particular, I will focus on identifying the emotional state of Charlie at the end of the dialogue, rather than simply relying on his verbal responses. I will consider how Charlie's tone, language, and behavior may indicate that he is feeling a particular emotion, such as embarrassment, provocation, indifference, or gratitude.} \\
\\ 
\texttt{To achieve this, I will analyze the dialogue more closely, looking for subtle cues and emotional undertones that may indicate Charlie's emotional state. I will also consider the implications of Charlie's words and actions, and how they may reflect his true feelings and intentions.} \\
\\ 
\texttt{With this revised approach, I will be better equipped to accurately predict Charlie's emotional response at the end of the dialogue.} \\
            }
        }
    }
    & \\ \\

\end{supertabular}
}

\end{document}
